% !TeX program = xelatex
% 第二章 2.6、2.7 純文字稿，接在既有 2.5 之後。
% 不含 chapter 或手動章節編號。

\section{任務流程與研究範圍}
\label{sec:task_workflow_scope}

本研究以室內取物任務作為系統整合情境，設定機器人由起點出發，經由視覺導引完成開門與通行，再夾取茶几上的圓柱體並持物返回。此任務串聯導航、目標接近、抓持與受約束操作，用以探討底盤及手臂在不同任務階段之間的協調方式。

\subsection{導航、接近與抓持流程}
\label{subsec:navigation_approach_grasping}

系統首先接收任務目的地、操作區域與目標物件類型等資訊，並依環境地圖建立導航參考。行進過程中，控制器根據底盤狀態及障礙資訊更新運動命令；當目標進入可觀測範圍時，視覺模組取得其位置資訊，供參考生成模組建立後續接近目標。

接近階段結合視覺量測、底盤位姿與手臂構型，調整底盤行進參考及末端執行器目標，使目標位於可操作範圍內。任務階段依末端到位、夾爪狀態及抓持條件推進，並持續檢查目標資料的有效性。流程設計以行進中觀測與底盤、手臂協同運動為主線，依當下狀態調整接近方向、速度與操作時機。

\subsection{開門、穿門與取物返回流程}
\label{subsec:door_traversal_retrieval_return}

當行進路徑經過關閉的門扇時，系統利用視覺量測取得把手位置，生成接近與抓持參考。抓持成立後，依門扇的旋轉機構及機器人狀態建立開門目標，協調底盤與手臂運動，使末端執行器的運動符合門扇機構限制。

門扇開啟至滿足通行需求後，機器人完成放手與手臂退離，再依門扇位置、通道幾何及障礙資訊更新導航參考，穿越門口並前往取物區域。通行過程需共同考量底盤占用範圍與手臂構型，使操作後的機器人姿態能銜接後續移動。

到達茶几附近後，系統透過視覺定位圓柱體，依量測結果及當前構型建立接近與抓取參考。抓持成立後，調整手臂至適合持物移動的構型，再由導航模組引導返回起點。返程持續考量障礙、門扇與持物構型對通行的影響。本研究之整合任務以持物返回為終點，暫不包含關門作業。

\subsection{研究範圍與驗證邊界}
\label{subsec:research_scope_validation_boundaries}

本研究以室內平面環境與指定類型物件為主要研究範圍。環境地圖、任務區域、物件類型及必要的幾何與機構資訊，依實驗設定提供；操作目標則由當次視覺量測與機器人狀態建立。研究重點為上述資訊條件下的預測導航、底盤與手臂協同控制，以及任務階段之間的銜接。

驗證以物理模擬為主，依動態導航、協同操作與視覺導引等情境分別評估，並以開門、穿門、取物及返回作為整合驗證目標。抽屜操作保留為關節式物件的協同控制比較案例。對採用理想抓持或模擬器接觸資訊的實驗，將明示其模型假設與資料用途；各項結果的適用範圍，依第四章所列實驗條件及實際完成情形界定。

\section{資料記錄與實驗重現方式}
\label{sec:data_recording_reproducibility}

為建立系統運作與實驗結果之間的對應關係，本研究以獨立實驗識別碼保存程式版本、場景配置、初始條件與控制參數。記錄內容包含底盤位姿、關節狀態、手臂位置設定點、操作目標、控制器輸出、實際放行命令及任務階段等資訊，使後續分析能追溯參考生成、命令處理與機器人運動之間的關係。

實驗資料以結構化檔案保存，並透過分析程式產生評估結果與圖表。模擬影像依時間與狀態記錄對應，用於檢查任務切換、接觸事件及失敗原因；模擬時間與主機實際運算時間則分別記錄，以分析運動表現與計算成本。各次實驗保留其配置及執行紀錄，作為後續重現與方法比較之依據，具體實驗條件與評估指標將於第四章說明。
